\chapter{Correspondência com o plano curricular}
\label{ap:curriculo}

Uma dissertação de mestrado em Engenharia de Inteligência Artificial deve poder mostrar onde é que
o conhecimento do curso foi efetivamente aplicado. Este apêndice faz essa correspondência, unidade
a unidade, e nomeia também as duas unidades cuja matéria não teve aplicação neste trabalho ---
porque uma correspondência que encontra tudo em todo o lado não informa ninguém.

\begin{table}[!htbp]
\caption[Correspondência com o plano curricular]{Onde cada unidade curricular do plano do mestrado
é aplicada neste trabalho. A última coluna remete para a secção onde a aplicação é visível.}
\label{tab:ap_curriculo}
\centering
\small
\begin{tabular}{L{0.24\textwidth} L{0.46\textwidth} L{0.20\textwidth}}
\toprule
\tabhead{Unidade} & \tabhead{Onde é aplicada} & \tabhead{Secção} \\
\midrule
\multicolumn{3}{l}{\emph{Fundamentos de sistemas inteligentes}} \\
Aprendizagem Automática 1 &
  Construção de um conjunto de dados rotulado a partir de duas fontes desalinhadas; regressão
  logística com normalização e reponderação de classes; conjuntos de árvores por reforço do
  gradiente como referência superior; calibração \emph{a posteriori} &
  §\ref{sec:met_triagem}, §\ref{sec:av_qi3} \\
Engenharia do Conhecimento &
  Raciocínio baseado em casos: recuperar e reutilizar precedentes em vez de aplicar uma regra
  geral, com a paragem deliberada antes da adaptação &
  §\ref{sec:ctx_casos}, §\ref{sec:met_recuperacao} \\
Paradigmas de Programação em IA &
  Separação entre lógica pura e efeitos, que é o que permite testar o núcleo sem rede e sem
  carregar modelos. A matéria de programação lógica e por restrições não teve aplicação &
  §\ref{sec:met_avaliacao} \\
Planeamento e Tomada de Decisão &
  Política de decisão sob orçamento: portas em cascata, limiar derivado de uma análise de custos
  de erro, e a distinção entre uma política \emph{offline} que conhece o dia inteiro e uma política
  \emph{online} que decide a cada ciclo &
  §\ref{sec:sis_portas}, §\ref{sec:av_qi3} \\
\midrule
\multicolumn{3}{l}{\emph{IA avançada e aplicações especializadas}} \\
Aprendizagem Automática 2 &
  Representações contextuais de frase produzidas por uma arquitetura \emph{Transformer}; exportação
  do codificador para um formato leve, com medição da fidelidade da conversão, por restrição de
  memória do ambiente de execução &
  §\ref{sec:ctx_texto}, §\ref{sec:met_recuperacao} \\
Linguagem Natural e Sistemas Conversacionais &
  Semelhança semântica entre frases e avaliação de recuperação sensível à ordem. O sistema é uma
  arquitetura de recuperação sem geração: recupera evidência e recusa produzir texto sobre ela, com
  a justificação dada no estado da arte &
  §\ref{sec:ctx_llm}, §\ref{sec:av_qi2} \\
Ambientes Inteligentes &
  Sem aplicação. O trabalho não tem componente física, robótica nem de Internet das Coisas &
  --- \\
\midrule
\multicolumn{3}{l}{\emph{Profissionalismo, segurança e inovação}} \\
Aspetos Sociais e Éticos em IA &
  Confiança apropriada em sistemas automáticos e o risco de a explicação aumentar a aceitação de
  respostas erradas; recusa de aconselhamento; fadiga de alertas como requisito de desenho &
  §\ref{sec:ctx_xai}, §\ref{sec:met_desafios} \\
Privacidade e Segurança em Sistemas Inteligentes &
  Minimização de dados por desenho, com a excepção declarada; gestão de credenciais e mascaramento
  de segredos em registos; verificação de integridade do modelo descarregado &
  §\ref{sec:met_privacidade}, §\ref{sec:met_seguranca} \\
Investigação e Inovação em IA &
  Enquadramento em investigação por desenho; critérios de sucesso fixados antes das medições;
  linhas de base declaradas com o respetivo chão; matriz de evidência com as afirmações retiradas;
  escrita científica e verificação das referências &
  §\ref{sec:met_tipo}, Apêndice~\ref{ap:reprodutibilidade} \\
\bottomrule
\end{tabular}
\end{table}

Duas observações sobre esta tabela.

A primeira é que a contribuição de engenharia deste trabalho não está no tamanho do modelo. É
treinado um único modelo, e é pequeno. Está no que foi necessário construir para que o número que
ele produz, e a conclusão negativa que esse número sustenta, sejam de confiança: o conjunto de
dados rotulado, a separação temporal imposta por um teste automático, a calibração com o
enviesamento declarado, a garantia de que o modelo implantado é o modelo avaliado, e o ciclo que
volta a avaliar o sistema depois de estar em funcionamento.

A segunda é que duas linhas dizem ``sem aplicação'' ou aplicação parcial. Ambientes Inteligentes
não tem correspondência neste trabalho, e a matéria de programação lógica e por restrições também
não. Registá-lo é mais informativo do que construir uma ligação artificial.
